\chapter{2000年广东卷}
\section{单选题}
\begin{problem}
已知集合 \(A = \{1, 2, 3, 4\}\)，那么 \(A\) 的真子集的个数是（\quad）

\choices
    {\(15\)}
    {\(16\)}
    {\(3\)}
    {\(4\)}
\end{problem}

\begin{answer}
A
\end{answer}

\begin{solution}
集合 \(A\) 有 \(4\) 个元素，因此它的全部子集共有
\[
2^{4}=16
\]
个. 真子集不包括集合 \(A\) 本身，所以真子集的个数为
\[
16-1=15
\]
故选 A.
\end{solution}

\begin{problem}
在复平面内，把复数 \(3 - \sqrt{3}\i\) 对应的向量按顺时针方向旋转 \(\frac{\pi}{3}\)，所得向量对应的复数是（\quad）

\choices
    {\(2\sqrt{3}\)}
    {\(-2\sqrt{3}\i\)}
    {\(\sqrt{3} - 3\i\)}
    {\(3 + \sqrt{3}\i\)}
\end{problem}

\begin{answer}
B
\end{answer}

\begin{solution}
顺时针旋转 \(\frac{\pi}{3}\) 相当于乘以
\[
\cos\left(-\frac{\pi}{3}\right)+\i\sin\left(-\frac{\pi}{3}\right)
=\frac12-\frac{\sqrt{3}}2\i
\]
因此所得复数为
\[
(3-\sqrt{3}\i)\left(\frac12-\frac{\sqrt{3}}2\i\right)
=\frac32-\frac{3\sqrt{3}}2\i-\frac{\sqrt{3}}2\i+\frac32\i^2
=-2\sqrt{3}\i
\]
故选 B.
\end{solution}

\begin{problem}
一个长方体共一顶点的三个面的面积分别是 \(\sqrt{2}\)，\(\sqrt{3}\)，\(\sqrt{6}\)，这个长方体对角线的长是（\quad）

\choices
    {\(2 \sqrt{3}\)}
    {\(3\sqrt{2}\)}
    {\(6\)}
    {\(\sqrt{6}\)}
\end{problem}

\begin{answer}
D
\end{answer}

\begin{solution}
设长方体从同一顶点出发的三条棱长分别为 \(x\)，\(y\)，\(z\)，
则 \(xy=\sqrt{2}\)，\(yz=\sqrt{3}\)，\(zx=\sqrt{6}\).
于是 \(x^{2}=\frac{(xy)(xz)}{yz}=2\)，\(y^{2}=\frac{(xy)(yz)}{xz}=1\)，\(z^{2}=\frac{(yz)(xz)}{xy}=3\).
长方体对角线长为
\[
\sqrt{x^{2}+y^{2}+z^{2}}=\sqrt{2+1+3}=\sqrt{6}
\]
故选 D.
\end{solution}

\begin{problem}
已知 \(\sin \alpha > \sin \beta\)，那么下列命题成立的是（\quad）

\choices
    {若 \(\alpha\)、\(\beta\) 是第一象限角，则 \(\cos \alpha > \cos \beta\)}
    {若 \(\alpha\)、\(\beta\) 是第二象限角，则 \(\tan \alpha > \tan \beta\)}
    {若 \(\alpha\)、\(\beta\) 是第三象限角，则 \(\cos \alpha > \cos \beta\)}
    {若 \(\alpha\)、\(\beta\) 是第四象限角，则 \(\tan \alpha > \tan \beta\)}
\end{problem}

\begin{answer}
D
\end{answer}

\begin{solution}
在每个象限内考察相应三角函数的单调性.

第一象限内，正弦递增而余弦递减，所以 \(\sin\alpha>\sin\beta\) 推出 \(\alpha>\beta\)，从而 \(\cos\alpha<\cos\beta\)，A 错.

第二象限内，正弦递减而正切递增，所以 \(\sin\alpha>\sin\beta\) 推出 \(\alpha<\beta\)，从而 \(\tan\alpha<\tan\beta\)，B 错.

第三象限内，正弦递减而余弦递增，所以 \(\sin\alpha>\sin\beta\) 推出 \(\alpha<\beta\)，从而 \(\cos\alpha<\cos\beta\)，C 错.

第四象限内，正弦递增且正切递增，所以 \(\sin\alpha>\sin\beta\) 推出 \(\alpha>\beta\)，进而 \(\tan\alpha>\tan\beta\). 故选 D.
\end{solution}

\begin{problem}
函数 \( y = -x \cos x \) 的部分图象是（\quad）

\choices
    {\choicebitmap[width=0.15\paperwidth]{img/2000/guangdong/q5_optA_nolabel.png}}
    {\choicebitmap[width=0.15\paperwidth]{img/2000/guangdong/q5_optB_nolabel.png}}
    {\choicebitmap[width=0.15\paperwidth]{img/2000/guangdong/q5_optC_nolabel.png}}
    {\choicebitmap[width=0.15\paperwidth]{img/2000/guangdong/q5_optD_nolabel.png}}
\end{problem}

\begin{answer}
D
\end{answer}

\begin{solution}
函数满足
\[
f(0)=0
\]
且 \(f'(x)=-\cos x+x\sin x\)，\(f'(0)=-1\).
所以图象在原点附近从左上方向右下方通过原点. 函数的零点满足
\[
-x\cos x=0
\]
即 \(x=0\)，或 \(x=\frac{\pi}{2}+k\pi\)（\(k\in\Z\)）. 特别地，在原点附近有
\[
f\left(-\frac{\pi}{2}\right)=f(0)=f\left(\frac{\pi}{2}\right)=0
\]
由 \(f'(0)=-1\)，曲线在原点处由左上方向右下方穿过. 再由零点的表达式可知，在 \((-\pi,-\frac{\pi}{2})\) 上 \(x<0\)、\(\cos x<0\)，故 \(f(x)<0\)；在 \((-\frac{\pi}{2},0)\) 上 \(x<0\)、\(\cos x>0\)，故 \(f(x)>0\). 在 \((0,\frac{\pi}{2})\) 上 \(x>0\)、\(\cos x>0\)，故 \(f(x)<0\)；在 \((\frac{\pi}{2},\pi)\) 上 \(x>0\)、\(\cos x<0\)，故 \(f(x)>0\). 因此曲线依次在 \(x=-\frac{\pi}{2}\)、\(x=0\)、\(x=\frac{\pi}{2}\) 处穿过横轴，穿越方向分别为向上、向下、向上. 逐项对照四个图形，只有第四个图形符合这些零点、符号和切线方向，故选 D.
\end{solution}

\begin{problem}
《中华人民共和国个人所得税法》规定，公民全月工资、薪金所得不超过 \(800\text{ 元}\) 的部分不必纳税. 超过 \(800\text{ 元}\) 的部分为全月应纳税所得额，此项税款按下列标准分段累进计算：

\begin{center}
\begin{tabular}{|c|c|}
\hline
全月应纳税所得额 & 税率 \\
\hline
不超过 \(500\text{ 元}\) 的部分 & \(5\%\) \\
\hline
超过 \(500\text{ 元}\) 至 \(2000\text{ 元}\) 的部分 & \(10\%\) \\
\hline
超过 \(2000\text{ 元}\) 至 \(5000\text{ 元}\) 的部分 & \(15\%\) \\
\hline
\(\cdots\) & \(\cdots\) \\
\hline
\end{tabular}
\end{center}

某人一月份交纳此项税款 \(26.78\text{ 元}\)，则他的当月工资、薪金所得介于（\quad）

\choices
    {\(800\text{ 元} - 900\text{ 元}\)}
    {\(900\text{ 元} - 1200\text{ 元}\)}
    {\(1200\text{ 元} - 1500\text{ 元}\)}
    {\(1500\text{ 元} - 2800\text{ 元}\)}
\end{problem}

\begin{answer}
C
\end{answer}

\begin{solution}
设月工资为 \(w\)，超过 \(800\) 元的部分为
\[
u=w-800
\]
前 \(500\) 元的税额为 \(500\times5\%=25\) 元. 由于总税额为 \(26.78\) 元，且
\[
25<26.78<25+(2000-500)\times10\%=175
\]
所以 \(500<u<2000\)，确实落在第二档. 此时税额为
\[
25+(u-500)\times10\%=0.1u-25
\]
令其等于 \(26.78\)，得 \(0.1u-25=26.78\)，\(u=517.8\)，\(w=1317.8\).
因为 \(1200<1317.8<1500\)，故选 C.
\end{solution}

\begin{problem}
若 \(a > b > 1\)，\(P = \sqrt{\lg a \cdot \lg b}\)，\(Q = \frac{1}{2} (\lg a + \lg b)\)，\(R = \lg \left(\frac{a + b}{2}\right)\)，则（\quad）

\choices
    {\(R < P < Q\)}
    {\(P < Q < R\)}
    {\(Q < P < R\)}
    {\(P < R < Q\)}
\end{problem}

\begin{answer}
B
\end{answer}

\begin{solution}
令 \(x=\lg a\)，\(y=\lg b\).
由 \(a>b>1\) 得 \(x>y>0\). 正数的几何平均数小于算术平均数，因此
\[
P=\sqrt{xy}<\frac{x+y}{2}=Q
\]
另一方面，
\[
\frac{a+b}{2}>\sqrt{ab}
\]
对两边取常用对数，得到
\[
R=\lg\left(\frac{a+b}{2}\right)>\lg\sqrt{ab}
=\frac{\lg a+\lg b}{2}=Q
\]
所以 \(P<Q<R\)，故选 B.
\end{solution}

\begin{problem}
以极坐标系中的点 \((1,1)\) 为圆心，\(1\) 为半径的圆的方程是（\quad）

\choices
    {\(\rho = 2\cos \left(\theta -\frac{\pi}{4}\right)\)}
    {\(\rho = 2\sin \left(\theta -\frac{\pi}{4}\right)\)}
    {\(\rho = 2\cos (\theta -1)\)}
    {\(\rho = 2\sin (\theta -1)\)}
\end{problem}

\begin{answer}
C
\end{answer}

\begin{solution}
极坐标点 \((1,1)\) 的直角坐标为
\[
(\cos 1,\sin 1)
\]
设圆上任一点的极坐标为 \((\rho,\theta)\). 由圆心距等于半径，得
\[
(\rho\cos\theta-\cos1)^2+(\rho\sin\theta-\sin1)^2=1
\]
展开并利用 \(\cos\theta\cos1+\sin\theta\sin1=\cos(\theta-1)\)，得
\[
\rho^2-2\rho\cos(\theta-1)=0
\]
对非原点的圆上点可约去 \(\rho\)，得到圆的标准极坐标方程
\[
\rho=2\cos(\theta-1)
\]
当 \(\rho=0\) 时表示的原点本身也在该圆上，并且在标准方程中由 \(\cos(\theta-1)=0\) 得到，因此不影响上述方程表示整圆. 故选 C.
\end{solution}

\begin{problem}
一个圆柱的侧面展开图是一个正方形，这个圆柱的全面积与侧面积的比例是（\quad）

\choices
    {\(\frac{1 + 2\pi}{2\pi}\)}
    {\(\frac{1 + 4\pi}{4\pi}\)}
    {\(\frac{1 + 2\pi}{\pi}\)}
    {\(\frac{1 + 4\pi}{2\pi}\)}
\end{problem}

\begin{answer}
A
\end{answer}

\begin{solution}
设圆柱底面半径为 \(r\)，高为 \(h\). 侧面展开为正方形，说明正方形的边长既是圆柱的高又是底面圆周长，因此
\[
h=2\pi r
\]
圆柱侧面积与全面积分别为 \(S_{\text{侧}}=2\pi rh\)，\(S_{\text{全}}=2\pi rh+2\pi r^2\).
所以
\[
\frac{S_{\text{全}}}{S_{\text{侧}}}
=1+\frac{r}{h}
=1+\frac{1}{2\pi}
=\frac{1+2\pi}{2\pi}
\]
故选 A.
\end{solution}

\begin{problem}
过原点的直线与圆 \(x^{2} + y^{2} + 4x + 3 = 0\) 相切，若切点在第三象限，则该直线的方程是（\quad）

\choices
    {\(y = \sqrt{3} x\)}
    {\(y = -\sqrt{3} x\)}
    {\( y = \frac{\sqrt{3}}{3} x \)}
    {\(y = -\frac{\sqrt{3}}{3} x\)}
\end{problem}

\begin{answer}
C
\end{answer}

\begin{solution}
圆的方程可写为
\[
(x+2)^2+y^2=1
\]
所以圆心为 \((-2,0)\)，半径为 \(1\). 设过原点的切线为 \(y=mx\)，即 \(mx-y=0\). 由圆心到切线的距离等于半径，得
\[
\frac{2\lvert m\rvert}{\sqrt{m^2+1}}=1
\]
平方并整理得 \(4m^2=m^2+1\)，所以 \(m^2=\frac13\).
所以 \(m=\pm\frac{\sqrt3}{3}\). 对直线 \(mx-y=0\)，圆心 \((-2,0)\) 到直线的垂足为
\[
\left(-\frac{2}{1+m^2},-\frac{2m}{1+m^2}\right)
\]
若切点在第三象限，垂足的两个坐标都应为负，因此取 \(m>0\). 因此 \(m=\frac{\sqrt{3}}3\)，直线方程为 \(y=\frac{\sqrt{3}}3x\).
故选 C.
\end{solution}

\begin{problem}
过抛物线 \(y = ax^{2}\)（\(a > 0\)）的焦点 \(F\) 作一直线交抛物线于 \(P\)，\(Q\) 两点，若线段 \(PF\) 与 \(FQ\) 的长分别是 \(p\)，\(q\)，则 \(\frac{1}{p} + \frac{1}{q}\) 等于（\quad）

\choices
    {\(2a\)}
    {\(\frac{1}{2a}\)}
    {\(4a\)}
    {\(\frac{4}{a}\)}
\end{problem}

\begin{answer}
C
\end{answer}

\begin{solution}
将抛物线写成标准形式
\[
x^2=\frac1a y=4p_0y
\]
其中 \(p_0=\frac{1}{4a}\)，
故焦点为 \(F=(0,p_0)\). 设过焦点的非竖直直线为
\[
y=mx+p_0
\]
它与抛物线的交点横坐标 \(x_1\)，\(x_2\) 满足
\[
ax^2-mx-p_0=0
\]
由韦达定理，
\[
x_1x_2=-\frac{p_0}{a}=-\frac{1}{4a^2}
\]
又
\[
|x_1-x_2|=\frac{\sqrt{m^2+4ap_0}}{a}=\frac{\sqrt{m^2+1}}a
\]
两根异号，所以
\[
|x_1|+|x_2|=|x_1-x_2|
\]
又因为点 \((x,mx+p_0)\) 到焦点的距离为 \(\lvert x\rvert\sqrt{1+m^2}\)，故
\[
p+q=(|x_1|+|x_2|)\sqrt{1+m^2}=\frac{1+m^2}{a}
\]
以及
\[
pq=|x_1x_2|(1+m^2)=\frac{1+m^2}{4a^2}
\]
因此
\[
\frac1p+\frac1q=\frac{p+q}{pq}=4a
\]
故选 C.
\end{solution}

\begin{problem}
如图，\(OA\) 是圆锥底面中心 \(O\) 到母线的垂线，\(OA\) 绕轴旋转一周所得曲面将圆锥分成体积相等的两部分，则母线与轴的夹角的余弦值为（\quad）

\bitmapfigure[max width=0.34\linewidth]{img/2000/guangdong/q12_fig1.png}

\choices
    {\(\frac{1}{\sqrt[3]{2}}\)}
    {\(\frac{1}{2}\)}
    {\(\frac{1}{\sqrt{2}}\)}
    {\(\frac{1}{\sqrt[4]{2}}\)}
\end{problem}

\begin{answer}
D
\end{answer}

\begin{solution}
设圆锥顶点为 \(B\)，底面圆心为 \(O\)，底面圆周上的点为 \(C\)，并令
\(OC=R\)，\(DA=r\)
其中 \(D\) 是过 \(A\) 向轴 \(OB\) 作垂线的垂足；设母线与轴的夹角为
\[
\angle OBC=\beta
\]
旋转轴截面中的三角形 \(OBA\) 时，在轴上以 \(D\) 为分界，它所生成的部分可分成两个底面半径均为 \(r\)、高分别为 \(OD\) 和 \(BD\) 的圆锥. 因此该部分体积为
\[
\frac13\pi r^2(OD+BD)=\frac13\pi r^2OB
\]
题设两部分体积相等，而整个圆锥体积为
\[
\frac13\pi R^2OB
\]
故
\(\frac13\pi r^2OB=\frac12\cdot\frac13\pi R^2OB\)，\(R^2=2r^2\).

由 \(OA\perp BC\)，线段 \(OA\) 在垂直于轴的方向上的投影为
\[
r=OA\cos\beta
\]
在直角三角形 \(OBC\) 与 \(OBA\) 中，分别有
\(R=OB\tan\beta\)，\(OA=OB\sin\beta\)
从而
\(OA=R\cos\beta\)，\(r=R\cos^2\beta\).
代入 \(R^2=2r^2\)，得
\(R^2=2R^2\cos^4\beta\)，\(\cos^4\beta=\frac12\).
由于 \(\beta\) 是锐角，
\[
\cos\beta=\frac1{\sqrt[4]{2}}
\]
故选 D.
\end{solution}

\section{填空题}
\begin{problem}
乒乓球队的 \(10\) 名队员中有 \(3\) 名主力队员，派 \(5\) 名参加比赛，\(3\) 名主力队员要安排在第一，三，五位置，其余 \(7\) 名队员选 \(2\) 名安排在第二，四位置，那么不同的出场安排共有\fillinblank{} 种（用数字作答）.
\end{problem}

\begin{answer}
252
\end{answer}

\begin{solution}
三名主力队员必须分别占据第一、三、五位置，有
\[
3!
\]
种安排. 其余两名队员从七名非主力队员中选出并安排在第二、四位置，有
\[
\mathrm C_7^2 2!=7\cdot6
\]
种安排. 因此总数为
\[
3!\mathrm C_7^2 2!=6\cdot7\cdot6=252
\]
\end{solution}

\begin{problem}
椭圆 \(\frac{x^{2}}{9} + \frac{y^{2}}{4} = 1\) 的焦点 \(F_{1}\)，\(F_{2}\)，点 \(P\) 为其上的动点，当 \(\angle F_{1}PF_{2}\) 为钝角时，点 \(P\) 横坐标的取值范围是\fillinblank{}.
\end{problem}

\begin{answer}
\(\left(-\frac{3}{\sqrt{5}},\frac{3}{\sqrt{5}}\right)\)
\end{answer}

\begin{solution}
椭圆的焦点为
\(F_1=(-\sqrt{5},0)\)，\(F_2=(\sqrt{5},0)\).
设动点 \(P=(x,y)\). 当 \(\angle F_1PF_2\) 为钝角时，两条向量的内积为负，即
\[
\overrightarrow{PF_1}\cdot\overrightarrow{PF_2}<0
\]
计算内积得
\[
(-x-\sqrt5,-y)\cdot(-x+\sqrt5,-y)=x^2+y^2-5<0
\]
即
\[
x^2+y^2<5
\]
由椭圆方程
\[
\frac{x^2}{9}+\frac{y^2}{4}=1
\]
得
\[
y^2=4\left(1-\frac{x^2}{9}\right)
\]
代入钝角条件，得到
\[
x^2+4\left(1-\frac{x^2}{9}\right)<5
\iff \frac59x^2<1
\iff x^2<\frac95
\]
所以
\[
-\frac3{\sqrt5}<x<\frac3{\sqrt5}
\]
\end{solution}

\begin{problem}
设 \(\{a_{n}\}\) 是首项为 \(1\) 的正项数列，且 \((n+1)a_{n+1}^{2} - na_{n}^{2} + a_{n+1}a_{n} = 0\)（\(n=1\)，\(2\)，\(3\)，\(\ldots\)），则它的通项公式是 \(a_{n}=\)\fillinblank{}.
\end{problem}

\begin{answer}
\(\frac{1}{n}\)
\end{answer}

\begin{solution}
由于数列各项为正，令
\[
r_n=\frac{a_{n+1}}{a_n}>0
\]
原递推式除以 \(a_n^2\)，得
\[
(n+1)r_n^2+r_n-n=0
\]
将其因式分解为
\[
\bigl((n+1)r_n-n\bigr)(r_n+1)=0
\]
因为 \(r_n>0\)，只能有
\(r_n=\frac{n}{n+1}\)，\(a_{n+1}=\frac{n}{n+1}a_n\).
由 \(a_1=1\)，逐项相乘得到
\[
a_n=a_1\prod_{k=1}^{n-1}\frac{k}{k+1}
=\frac1n
\]
\end{solution}

\begin{problem}
如图，\(E\)，\(F\) 分别为正方体面 \(ADD_{1}A_{1}\)，面 \(BCC_{1}B_{1}\) 的中心，则四边形 \(BFD_{1}E\) 在该正方体的面上的射影可能是\fillinblank{}.

（要求：把可能的图序号都填上）

\bitmapfigure[max width=0.82\linewidth]{img/2000/guangdong/q16_fig1.png}
\end{problem}

\begin{answer}
\(\circled{2}\)，\(\circled{3}\)
\end{answer}

\begin{solution}
建立正方体坐标系，取
\(A=(0,0,0)\)，\(B=(1,0,0)\)，\(C=(1,1,0)\)，\(D=(0,1,0)\)
顶面各点的第三坐标为 \(1\). 则
\(E=\left(0,\frac12,\frac12\right)\)，\(F=\left(1,\frac12,\frac12\right)\)，\(D_1=(0,1,1)\).

向 \(x=0\) 的面作正投影时，\(B\) 投影到 \(A\)，\(F\) 与 \(E\) 的投影重合，\(D_1\) 不变；三点 \(A\)，\(E\)，\(D_1\) 共线. 向 \(x=1\) 的面作正投影时，\(A\) 投影到 \(B\)，\(E\) 与 \(F\) 的投影重合，\(D_1\) 投影到 \(C_1\)；三点 \(B\)，\(F\)，\(C_1\) 共线. 这两种情形均得到题图中的 \(\circled{3}\).

向 \(y=0\)、\(y=1\)、\(z=0\) 或 \(z=1\) 的面作正投影时，按 \(B\)，\(F\)，\(D_1\)，\(E\) 的顺序，四个投影点的两组对边分别平行且等长，所得图形为题图中的 \(\circled{2}\). 因此其余四个面的投影均为平行四边形，不会得到图形 \(\circled{1}\) 或 \(\circled{4}\).

所以可能的图序号为 \(\circled{2}\)，\(\circled{3}\).
\end{solution}

\section{解答题}
\begin{problem}
已知函数 \(y = \sqrt{3}\sin x + \cos x\)，\(x \in \R\).

\begin{enumerate}
\item 当函数 \(y\) 取得最大值时，求自变量 \(x\) 的集合；
\item 该函数的图象可由 \(y = \sin x\)（\(x \in \R\)）的图象经过怎样的平移和伸缩变换得到？
\end{enumerate}
\end{problem}

\begin{solution}
\begin{enumerate}
\item 将函数改写为
\[
y=\sqrt{3}\sin x+\cos x
=2\left(\frac{\sqrt{3}}2\sin x+\frac12\cos x\right)
=2\sin\left(x+\frac{\pi}{6}\right)
\]
因此 \(y\) 的最大值为 \(2\). 取最大值的充要条件是
\[
\sin\left(x+\frac{\pi}{6}\right)=1
\]
所以
\(x+\frac{\pi}{6}=\frac{\pi}{2}+2k\pi\)，\(k\in\Z\)
即
\(x=\frac{\pi}{3}+2k\pi\)，\(k\in\Z\).
\item 从 \(y=\sin x\) 的图象出发，先向左平移 \(\frac{\pi}{6}\)，得到
\[
y=\sin\left(x+\frac{\pi}{6}\right)
\]
的图象；再把纵坐标伸长到原来的 \(2\) 倍，得到
\[
y=2\sin\left(x+\frac{\pi}{6}\right)
=\sqrt{3}\sin x+\cos x
\]
\end{enumerate}
\end{solution}

\begin{problem}
设 \(\{a_{n}\}\) 为等比数列，\(T_{n} = na_{1} + (n - 1)a_{2} + \cdots + 2a_{n-1} + a_{n}\)，已知 \(T_{1} = 1\)，\(T_{2} = 4\).

\begin{enumerate}
\item 求数列 \(\{a_{n}\}\) 的首项和公比；
\item 求数列 \(\{T_{n}\}\) 的通项公式.
\end{enumerate}
\end{problem}

\begin{solution}
\begin{enumerate}
\item 由 \(T_1=a_1=1\)，得
\[
a_1=1
\]
又
\[
T_2=2a_1+a_2=4
\]
所以
\[
a_2=4-2a_1=2
\]
等比数列的公比为
\[
q=\frac{a_2}{a_1}=2
\]
从而
\[
a_n=a_1q^{n-1}=2^{n-1}
\]
\item 由定义直接相减，得
\[
T_n-T_{n-1}=a_1+a_2+\cdots+a_n
\]
由等比数列通项公式，
\[
a_1+a_2+\cdots+a_n=\sum_{k=0}^{n-1}2^k=2^n-1
\]
因此
\[
T_n=T_1+\sum_{j=2}^{n}(T_j-T_{j-1})
=1+\sum_{j=2}^{n}(2^j-1)
=1+(2^{n+1}-4)-(n-1)
=2^{n+1}-n-2
\]
\end{enumerate}
\end{solution}

\begin{problem}
如图，已知平行六面体 \(ABCD - A_{1}B_{1}C_{1}D_{1}\) 的底面 \(ABCD\) 是菱形，且 \(\angle C_{1}CB = \angle C_{1}CD = \angle BCD\).

\begin{enumerate}
\item 证明：\(C_{1}C \perp BD\)；
\item 当 \(\frac{CD}{CC_{1}}\) 的值为多少时，能使 \(A_{1}C \perp\) 平面 \(C_{1}BD\)？请给出证明.
\end{enumerate}

\bitmapfigure[max width=0.43\linewidth]{img/2000/guangdong/q19_fig1.png}
\end{problem}

\begin{solution}
设
\(\bs{b}=\overrightarrow{CB}\)，\(\bs{d}=\overrightarrow{CD}\)，\(\bs{u}=\overrightarrow{CC_1}\)
并记
\(|\bs{b}|=|\bs{d}|=s\)，\(|\bs{u}|=h\).
由于底面 \(ABCD\) 是菱形，且三个给定角相等，设它们均为 \(\theta\)，则
\[
\bs{u}\cdot\bs{b}=hs\cos\theta
=\bs{u}\cdot\bs{d}
\]
所以
\[
\bs{u}\cdot(\bs{d}-\bs{b})=0
\]
即
\[
CC_1\perp BD
\]

取 \(C\) 为向量起点，则
\[
\overrightarrow{CA_1}=\bs{b}+\bs{d}+\bs{u}
\]
平面 \(C_1BD\) 内的两个相交方向向量可取
\(\overrightarrow{BD}=\bs{d}-\bs{b}\)，\(\overrightarrow{BC_1}=\bs{u}-\bs{b}\).
先有
\[
(\bs{b}+\bs{d}+\bs{u})\cdot(\bs{d}-\bs{b})=0
\]
因为 \(|\bs{b}|=|\bs{d}|\) 且 \(\bs{u}\cdot(\bs{d}-\bs{b})=0\). 因此只需令
\[
(\bs{b}+\bs{d}+\bs{u})\cdot(\bs{u}-\bs{b})=0
\]
由 \(\bs{b}\cdot\bs{d}=s^2\cos\theta\) 及上面的内积关系，左边为
\[
h^2+hs\cos\theta-s^2(1+\cos\theta)
=(h-s)\bigl(h+s(1+\cos\theta)\bigr)
\]
由于 \(h\)，\(s>0\) 且 \(1+\cos\theta>0\)，第二因子为正，所以该内积为零当且仅当
\[
h=s
\]
于是
\[
\frac{CD}{CC_1}=\frac{s}{h}=1
\]
\end{solution}

\begin{problem}
设函数 \(f(x) = \sqrt{x^2 + 1} - ax\)，其中 \(a > 0\).

\begin{enumerate}
\item 解不等式 \(f(x) \leqslant 1\)；
\item 证明：当 \(a \geqslant 1\) 时，函数 \(f(x)\) 在区间 \(\left[0,+\infty\right)\) 上是单调函数.
\end{enumerate}
\end{problem}

\begin{solution}
\begin{enumerate}
\item 当 \(x<0\) 时，
\(\sqrt{x^2+1}>1\)，\(-ax>0\)
所以 \(f(x)>1\)，不等式无负数解. 令 \(x\geqslant0\)，则
\[
\sqrt{x^2+1}\leqslant1+ax
\]
且右边为正，可以平方而不改变不等号方向，得到
\[
x^2+1\leqslant(1+ax)^2
\iff x\bigl((1-a^2)x-2a\bigr)\leqslant0
\]
当 \(0<a<1\) 时，因 \(1-a^2>0\)，故
\[
0\leqslant x\leqslant\frac{2a}{1-a^2}
\]
当 \(a\geqslant1\) 时，对任意 \(x\geqslant0\)，有
\[
(1-a^2)x-2a<0
\]
所以不等式恒成立. 综上，解集为
\[
\begin{cases}
\left[0,\dfrac{2a}{1-a^2}\right],&0<a<1,\\[6pt]
\left[0,+\infty\right),&a\geqslant1.
\end{cases}
\]
\item 在 \([0,+\infty)\) 上，
\[
f'(x)=\frac{x}{\sqrt{x^2+1}}-a
\]
由于
\[
0\leqslant\frac{x}{\sqrt{x^2+1}}<1\leqslant a
\]
故 \(f'(x)<0\). 因此 \(f(x)\) 在 \([0,+\infty)\) 上严格单调递减.
\end{enumerate}
\end{solution}

\begin{problem}
某蔬菜基地种植西红柿，由历年市场行情得知，从二月一日起的 \(300\) 天内，西红柿市场售价与上市时间的关系用图一的一条折线表示；西红柿的种植成本与上市时间的关系用图二的抛物线段表示.

\begin{enumerate}
\item 写出图一表示的市场售价与时间的函数关系式 \(P = f(t)\) 和图二表示的种植成本与时间的函数关系式 \(Q = g(t)\)；
\item 认定市场售价减去种植成本为纯收益，问何时上市的西红柿纯收益最大？
\end{enumerate}

注：市场售价与种植成本的单位：\(\text{元}/10^{2}\mathrm{~kg}\)，时间单位：\(\text{天}\).
\FigureLayoutDeclare{img/2000/guangdong/q21_fig1.png}{11.7cm}{10bp 9bp 9bp 10bp}

\bitmapfigure[max width=0.75\linewidth]{img/2000/guangdong/q21_fig1.png}
\end{problem}

\begin{solution}
由图一两段直线分别经过 \((0,300)\)，\((200,100)\) 和 \((200,100)\)，\((300,300)\)，故
\[
f(t)=
\begin{cases}
300-t,&0\leqslant t\leqslant200,\\
2t-300,&200<t\leqslant300.
\end{cases}
\]
图二的抛物线顶点为 \((150,100)\)，且经过 \((50,150)\)，所以
\(g(t)=\frac1{200}(t-150)^2+100\)，\(0\leqslant t\leqslant300\).
设纯收益为 \(h(t)=f(t)-g(t)\). 当 \(0\leqslant t\leqslant200\) 时，
\[
h(t)=300-t-\frac1{200}(t-150)^2-100
=100-\frac{(t-50)^2}{200}
\]
故这一段的最大值为 \(100\)，在 \(t=50\) 时取得.

当 \(200<t\leqslant300\) 时，
\[
h'(t)=2-\frac{t-150}{100}=\frac{350-t}{100}>0
\]
所以 \(h(t)\) 在这一段递增，最大值在 \(t=300\) 时取得，且
\[
h(300)=300-\left(\frac{150^2}{200}+100\right)=87.5
\]
由于 \(100>87.5\)，全区间最大收益在 \(t=50\) 时取得，即从二月一日开始的第 \(50\) 天.
\end{solution}

\begin{problem}
如图，已知梯形 \(ABCD\) 中 \(|AB| = 2|CD|\)，点 \(E\) 分有向线段 \(\overrightarrow{AC}\) 所成的比为 \(\lambda\)，双曲线过 \(C\)，\(D\)，\(E\) 三点，且以 \(A\)，\(B\) 为焦点，当 \(\frac{2}{3} \leqslant \lambda \leqslant \frac{3}{4}\) 时，求双曲线离心率 \(e\) 的取值范围.
\bitmapfigure[max width=0.58\linewidth]{img/2000/guangdong/q22_fig1.png}
\end{problem}

\begin{solution}
建立直角坐标系，使 \(AB\) 为 \(x\) 轴，\(AB\) 的垂直平分线为 \(y\) 轴. 令
\(AB=2c\)，\(CD=c\)
则可设
\(A=(-c,0)\)，\(B=(c,0)\)，\(C=\left(\frac c2,h\right)\)，\(D=\left(-\frac c2,h\right)\)
其中 \(h>0\) 为梯形的高. 这里利用了双曲线以 \(A\)，\(B\) 为焦点，故其关于 \(y\) 轴对称，且 \(C\)，\(D\) 的纵坐标相同.

设双曲线的标准方程为
\(\frac{x^2}{a^2}-\frac{y^2}{b^2}=1\)，\(b^2=c^2-a^2\)
其离心率为 \(e=c/a\). 令
\(t=\frac{a^2}{c^2}\)，\(u=\frac hc\).
将点 \(C\) 代入双曲线方程，得
\[
\frac{1}{4t}-\frac{u^2}{1-t}=1
\]
即
\[
\frac{u^2}{1-t}=\frac{1-4t}{4t}
\]
由定比分点公式以及题设有向线段比，点 \(E\) 的坐标为
\[
E=\frac{A+\lambda C}{1+\lambda}
=\left(\frac{c(\lambda-2)}{2(1+\lambda)},
\frac{\lambda h}{1+\lambda}\right)
\]
将 \(E\) 代入双曲线方程，并除以相应的 \(c^2\)，得到
\[
\frac{(\lambda-2)^2}{4(1+\lambda)^2t}
-\frac{\lambda^2u^2}{(1+\lambda)^2(1-t)}=1
\]
利用
\[
\frac{u^2}{1-t}=\frac{1-4t}{4t}
\]
上式化为
\[
\frac{(\lambda-2)^2-\lambda^2(1-4t)}
{4(1+\lambda)^2t}=1
\]
整理分子后，得到
\[
\frac{1-\lambda+\lambda^2t}{(1+\lambda)^2t}=1
\]
从而
\(1-\lambda=(1+2\lambda)t\)，\(t=\frac{1-\lambda}{1+2\lambda}\).
因此
\[
e^2=\frac{c^2}{a^2}
=\frac1t
=\frac{1+2\lambda}{1-\lambda}
\]
函数
\[
\frac{1+2\lambda}{1-\lambda}
\]
在 \(\lambda<1\) 时严格递增，因为其导数为
\[
\frac{3}{(1-\lambda)^2}>0
\]
所以当 \(\frac23\leqslant\lambda\leqslant\frac34\) 时，
\[
7=\frac{1+2\cdot\frac23}{1-\frac23}
\leqslant e^2
\leqslant
\frac{1+2\cdot\frac34}{1-\frac34}=10
\]
又 \(e>0\)，故
\[
e\in\left[\sqrt7,\sqrt{10}\right]
\]
\end{solution}
